% Figure: light-current (P-I) curve of a Fabry-Perot laser diode against a
% bare LED. Below threshold the laser emits weak spontaneous light like an
% LED; above threshold stimulated emission switches on and the output rises
% steeply and linearly with a slope set by the differential quantum efficiency.
% Two laser curves at two temperatures show the threshold shift with heat.
% Sampled coordinates only, no computed sqrt.
\begin{figure}[htbp]
\centering
\begin{tikzpicture}[font=\scriptsize]
\begin{axis}[
    width=11cm, height=6.6cm,
    xlabel={drive current $I$ (\si{\milli\ampere})},
    ylabel={optical output power $P$ (\si{\milli\watt})},
    xmin=0, xmax=80,
    ymin=0, ymax=14,
    xtick={0,10,20,30,40,50,60,70,80},
    ytick={0,2,4,6,8,10,12,14},
    grid=both,
    grid style={enggray!25},
    tick label style={font=\tiny},
    label style={font=\scriptsize},
    axis line style={enggray},
    legend style={font=\tiny, at={(0.02,0.98)}, anchor=north west,
      draw=enggray, fill=white},
    every axis plot/.append style={very thick},
]
% laser at 25 C: threshold near 20 mA, slope efficiency ~0.3 mW/mA above it
\addplot[engblue] coordinates {
  (0,0.00) (5,0.02) (10,0.05) (15,0.09) (18,0.14) (20,0.20)
  (22,0.80) (25,1.60) (30,3.10) (40,6.10) (50,9.10) (60,12.1) (70,15.1)
};
\addlegendentry{laser, $25\,{}^{\circ}\mathrm{C}$}
% laser at 60 C: threshold shifted up to ~30 mA, slope slightly lower
\addplot[engred, dashed] coordinates {
  (0,0.00) (5,0.02) (10,0.04) (20,0.10) (25,0.16) (28,0.22) (30,0.30)
  (33,0.85) (36,1.60) (42,3.10) (52,5.70) (62,8.30) (72,10.9)
};
\addlegendentry{laser, $60\,{}^{\circ}\mathrm{C}$}
% LED-like spontaneous output: gentle sub-linear, no threshold
\addplot[engorange, densely dotted] coordinates {
  (0,0) (10,0.30) (20,0.55) (30,0.78) (40,0.98) (50,1.16)
  (60,1.32) (70,1.46) (80,1.58)
};
\addlegendentry{bare LED (spontaneous)}
% threshold markers
\addplot[engblue, only marks, mark=o, mark size=2.5pt] coordinates {(20,0.20)};
\node[engblue, anchor=south east, font=\tiny] at (axis cs:20,0.6)
  {$I_{th}\approx 20$};
\addplot[engred, only marks, mark=o, mark size=2.5pt] coordinates {(30,0.30)};
\node[engred, anchor=south west, font=\tiny] at (axis cs:31,0.4)
  {$I_{th}\approx 30$};
% slope annotation
\node[engblue, anchor=west, font=\tiny] at (axis cs:48,11.0)
  {slope $\eta_d\approx 0.3$ mW/mA};
\end{axis}
\end{tikzpicture}
\caption[Light-current curve of a laser diode]{Output power against drive
current for a Fabry-Perot laser diode at two heat-sink temperatures, with the
spontaneous output of a bare LED for contrast. Below the threshold current the
laser emits only weak spontaneous light and behaves like a dim LED; above
threshold the round-trip gain equals the cavity loss, stimulated emission
switches on, and the power rises steeply and almost linearly with a slope set
by the differential (slope) efficiency. Heating the device raises the threshold
current and lowers the slope, the temperature sensitivity captured by the
characteristic temperature $T_0$ and the reason laser transmitters carry a
thermoelectric cooler or a monitor-photodiode feedback loop.}
\label{fig:vol04ch12:laser-pi-curve}
\end{figure}
